\chapter{Mixed Quantum-Classical Dynamics}\label{sec:mixed_quantum_classical_dynamics}
\section{Quantum Amplitude Propagation in Mixed Schemes}\label{sec:quantum_amplitude_propagation_in_mixed_themes}
Under the \acrfull{boa}, the total molecular wavefunction, $\ket{\Psi(t)}$, is separable into electronic and nuclear parts because the nuclear masses greatly exceed those of the electrons. However, when two or more electronic states approach degeneracy, the \acrshort{boa} separation fails and nonadiabatic effects become significant. \acrfull{tsh} is a general framework that treats nuclei and electrons separately: nuclei are propagated classically on an adiabatic \acrfull{pes}, while electrons are treated quantum mechanically.

In \acrshort{tsh}, the nuclei follow Newton's equations of motion on a single adiabatic \acrshort{pes} at any instant:
\begin{equation}
\symbf{M}\ddot{\symbf{R}}(t)=-\nabla_{\symbf{R}}E_I,
\end{equation}
where $\symbf{M}$ is the diagonal matrix with masses for each coordinate, $\symbf{R}(t)$ is nuclear geometry at time $t$, and $E_I$ is the electronic energy of the state $I$ at the nuclear configuration $\symbf{R}(t)$. Meanwhile, the electrons evolve according to the \acrshort{tdse}, with the electronic wavefunction $\ket{\Phi(t)}$ expanded in the instantaneous adiabatic eigenvectors of the electronic Hamiltonian $H_{\symrm{el}}$ as
\begin{equation}
\ket{\Phi(t)}=\sum_I c_I(t)\ket{\phi_I},
\end{equation}
where $\phi_I$ satisfies the \acrfull{tise}
\begin{equation}
H_{\symrm{el}}\ket{\phi_I}=E_I\ket{\phi_I},
\end{equation}
and the complex coefficients $c_I$ encode the instantaneous probability amplitudes for occupying each adiabatic state. For notational brevity, we will often suppress explicit dependence on the variables where no ambiguity arises.

Substituting this ansatz into the full electronic \acrshort{tdse},
\begin{equation}
i\frac{\partial}{\partial t}\ket{\Phi(t)}=H_{\symrm{el}}\ket{\Phi(t)},
\end{equation}
and projecting onto a particular adiabatic state $\ket{\phi_I}$, we obtain a set of coupled equations for the coefficients $c_I$:
\begin{equation}\label{eq:tsh_eom}
i\frac{\symrm{d}c_I}{\symrm{d}t}=\sum_J\left(H_{IJ}^{\symrm{el}}-i\Braket{\phi_I\vert\frac{\partial\phi_J}{\partial t}}\right)c_J,
\end{equation}
where $\sigma_{IJ}=\Braket{\phi_I\vert\frac{\partial\phi_J}{\partial t}}$ is the \acrfull{tdc}. In the purely adiabatic basis, $H_{\symrm{el}}$ is diagonal, so electronic population transfer between states is mediated entirely by $\sigma_{IJ}$:
\begin{equation}\label{eq:sh_tdse_coefs}
i\dot c_I=E_I c_I-i\sum_J\sigma_{IJ}c_J.
\end{equation}
With $\symbf{c}(t)=\left(c_1(t),c_2(t),\dots,c_N(t)\right)^{\symrm{T}}$ the Eq~\eqref{eq:sh_tdse_coefs} is a Schr\"odinger equation in a finite-dimensional Hilbert space $\mathbb{C}^N$ and can be rewritten as
\begin{equation}
i\symbf{\dot{c}}(t)=H_{\symrm{eff}}\symbf{c}(t),\qquad H_{IJ}^{\symrm{eff}}=E_I\delta_{IJ}-i\sigma_{IJ}
\end{equation}
For electronic dynamics in Eq~\eqref{eq:sh_tdse_coefs} to be unitary, the effective Hamiltonian $H_{\symrm{eff}}$ must be self-adjoint, which in turn requires $\sigma_{IJ}$ to be antisymmetric, $\sigma^\dagger=-\sigma$, in accordance with Sec~\ref{sec:self_adjoint_operators_and_observables}. To evaluate $\sigma_{IJ}$ in practice, note that each adiabatic eigenfunction depends parametrically on the instantaneous nuclear coordinates $\symbf{R}(t)$. By the chain rule,
\begin{equation}
\sigma_{IJ}=\Braket{\phi_I\vert\frac{\partial\phi_J}{\partial t}}=\dot{\symbf{R}}\cdot\Braket{\phi_I\vert\frac{\partial\phi_J}{\partial\symbf{R}}}=\symbf{v}\cdot\symbf{d}_{IJ},
\end{equation}
where $\symbf{v}=\dot{\symbf{R}}$ is the nuclear velocity and $\symbf{d}_{IJ}$ is the \acrfull{nacv}. Although many quantum‐chemistry packages provide \acrshort{nacv} and thus allow computation of the \acrshort{tdc} via $\sigma_{IJ}=\symbf{v}\cdot\symbf{d}_{IJ}$, \acrshort{nacv} can become ill-defined or numerically unstable near conical intersections (e.g., in multi-reference methods). In cases where a reliable \acrshort{nacv} is not available, one may instead employ a direct \acrshort{tdc} approximation using wavefunction overlaps.
\section{The Time Derivative Coupling}\label{sec:time_derivative_coupling}
\acrshort{tdc} is a pivotal quantity in \acrshort{tsh} simulations because it enters the electronic equations of motion given in Eq.~\eqref{eq:tsh_eom}. In \textit{ab initio} molecular dynamics, the \acrshort{tdc} is commonly evaluated using the full \acrshort{nacv} as $\sigma_{IJ}=\symbf{v}\cdot\symbf{d}_{IJ}$, where $\symbf{v}$ is the nuclear velocity. When the dynamics evolves on analytic \acrshort{pes} or the \acrshort{nacv} is for any reason not available, it is often preferable to calculate the \acrshort{tdc} directly. The most straightforward strategy applies a finite-difference formula to the adiabatic wavefunctions
\begin{equation}\label{eq:tdc_fd1}
\sigma_{IJ}(t)=\frac{1}{\Delta t}\left(\Braket{\phi_I(t)\vert\phi_J(t+\Delta t)}-\Braket{\phi_I(t)\vert\phi_J(t)}\right)=\frac{1}{\Delta t}\left(\Braket{\phi_I(t)\vert\phi_J(t+\Delta t)}-\delta_{IJ}\right),
\end{equation}
with $\Delta t$ a small time step. When analytic \acrshort{pes} is available, the finite-difference estimate of the \acrshort{tdc} is straightforward to implement. In \textit{ab initio} molecular dynamics, however, it is usually more practical to recover this quantity from the non-adiabatic coupling vector, which most electronic-structure packages provide natively. The naïve first-order finite-difference formula suffers from well-known deficiencies: its truncation error scales linearly with the time step and it is highly sensitive to numerical noise. A simple remedy is to adopt the second-order approximation
\begin{equation}\label{eq:tdc_fd2}
\sigma_{IJ}(t)=\frac{1}{2\Delta t}\left(\Braket{\phi_I(t)\vert\phi_J(t+\Delta t)}-\Braket{\phi_I(t)\vert\phi_J(t-\Delta t)}\right).
\end{equation}
Although second-order accurate, the approximation still retains significant shortcomings. In particular, it still violates the fundamental antisymmetry requirement, $\sigma_{IJ}^\ast\neq-\sigma_{JI}$, undermining the Hermiticity of the electronic Hamiltonian. The following sections tackle these deficiencies and introduce more accurate, numerically robust strategies for evaluating the \acrshort{tdc}.
\subsection{The Hammes-Schiffer Tully Scheme}\label{sec:hammes_schiffer_tully_scheme}
The \acrfull{hst}\supercite{10.1063/1.467455} scheme remedies the antisymmetry defect of Eq.~\eqref{eq:tdc_fd1} and Eq.~\eqref{eq:tdc_fd2}, making the electronic Hamiltonian Hermitian. To that end we first introduce a midpoint approximation for each adiabatic eigenstate $\ket{\phi_I}$, obtained by linear interpolation between its values at the bracketing time slices
\begin{equation}\label{eq:hst_wfn}
\Ket{\phi_I\left(t+\frac{\Delta t}{2}\right)}=\frac{1}{2}\left(\Ket{\phi_I(t+\Delta t)}+\Ket{\phi_I(t)}\right).
\end{equation}
This midpoint state is not strictly normalized, a fully norm-preserving variant will be introduced in the next section. The time derivative at the same midpoint is obtained by combining a backward difference for $\phi_I(t+\Delta t)$ with a forward difference for $\phi_I(t)$, yielding
\begin{equation}\label{eq:hst_dwfn}
\Ket{\frac{\partial \phi_I}{\partial t}\left(t+\frac{\Delta t}{2}\right)}=\frac{1}{\Delta t}\left(\Ket{\phi_I(t+\Delta t)}-\Ket{\phi_I(t)}\right).
\end{equation}
The \acrshort{tdc} $\sigma_{IJ}=\Braket{\phi_I\vert\frac{\partial \phi_J}{\partial t}}$ at a time $t+\frac{\Delta t}{2}$ is then formed as an overlap beween the expressions in Eq.~\eqref{eq:hst_wfn} and Eq.~\eqref{eq:hst_dwfn}:
\begin{equation}\label{eq:hst_tdc}
\sigma_{IJ}\left(t+\frac{\Delta t}{2}\right)=\frac{1}{2\Delta t}\left(\Braket{\phi_I(t)\vert\phi_J(t+\Delta t)}-\Braket{\phi_I(t+\Delta t)\vert\phi_J(t)}\right).
\end{equation}
The \acrshort{hst} scheme in Eq.~\eqref{eq:hst_tdc} is a finite-difference approximation of the \acrshort{tdc}, which is also antisymmetric, i.e., $\sigma_{IJ}^\ast=-\sigma_{JI}$. The scheme is numerically stable, but it does not preserve the norm of the adiabatic wavefunction in between the timesteps. Any rapidly varying wavefunction will suffer from a significant normalization error, which can lead to unphysical results. Such a case occurs for two diabatic surfaces coupled by zero diabatic interaction: in the adiabatic basis the exact \acrshort{tdc} diverges, yet the \acrshort{hst} formula collapses to zero, predicting no population transfer. To address this issue, we can apply a \acrfull{npi} to the adiabatic wavefunction, which is discussed in the next section.
\subsection{The Norm-Preserving Interpolation}\label{sec:norm_preserving_interpolation}
The \acrshort{npi} method improves upon the \acrshort{hst} scheme by ensuring that the interpolated adiabatic wavefunctions remain normalized at all times.\supercite{10.1021/jz5009449,10.1021/acs.jctc.6b00673} This might seem like a minor technical detail, but it has significant implications for the accuracy and stability of the \acrshort{tdc} evaluation, especially in scenarios involving rapidly varying wavefunctions or near-degenerate states.

We start with the two-state case for simplicity. Assume we have two normalized adiabatic states $\ket{\phi_1(t)}$ and $\ket{\phi_2(t)}$ from the Hilbert space $\mathcal{H}$ at time $t$. Our integration time step is $\Delta t$. We can write all the intermediate states between $t$ and $t+\Delta t$ as
\begin{equation}\label{eq:npi_interpolation}
\ket{\phi_I\left(\tau\right)}=U(\tau)\ket{\phi_I(t)},
\end{equation}
where $\tau\in[t,t+\Delta t]$ and $U(\tau)$ is a unitary operator that smoothly transforms the states from time $t$ to time $\tau$. We want to construct $U(\tau)$ such that it preserves the normalization of the states at all times, i.e., $\Braket{\phi_I(\tau)\vert\phi_I(\tau)}=1$ for all $\tau$. A simple choice for $U(\tau)$ is a rotation in the two-dimensional subspace spanned by $\ket{\phi_1(t)}$ and $\ket{\phi_2(t)}$ defined from the well-known 2D rotation matrix as
\begin{equation}\label{eq:npi_unitary_transform}
U_{II}(\tau)=\cos\left(\frac{\arccos\left(U_{II}\left(t+\Delta t\right)\right)}{\Delta t}\left(\tau-t\right)\right),\quad U_{IJ}(\tau)=\sin\left(\frac{\arcsin\left(U_{IJ}\left(t+\Delta t\right)\right)}{\Delta t}\left(\tau-t\right)\right),
\end{equation}
where
\begin{equation}\label{eq:npi_unitary_matrix}
U_{IJ}\left(t+\Delta t\right)=\Braket{\phi_I(t)\vert\phi_J(t+\Delta t)}.
\end{equation}
We could easily verify that the matrix $U(\tau)$ is unitary for all $\tau\in[t,t+\Delta t]$ and thus preserves the normalization of the states.
Using the interpolation in Eq~\eqref{eq:npi_interpolation} and the unitary transformation in Eq~\eqref{eq:npi_unitary_transform}, we can now compute the \acrshort{tdc} in arbitrary time $\tau\in[t,t+\Delta t]$ as
\begin{equation}\label{eq:npi_overlap}
\sigma_{IJ}(\tau)=\Braket{\phi_I(t)\vert U^\dagger(\tau)\frac{\partial}{\partial\tau}U(\tau)\vert\phi_J(t)}.
\end{equation}
Averaging over the interval $[t,t+\Delta t]$ gives the final \acrshort{npi} expression for the \acrshort{tdc}
\begin{equation}\label{eq:npi_tdc_final}
\sigma_{IJ}=\frac{1}{\Delta t}\int_t^{t+\Delta t}\Braket{\phi_I(t)\vert U^\dagger(\tau)\frac{\partial}{\partial\tau}U(\tau)\vert\phi_J(t)}\symrm{d}\tau.
\end{equation}
We explicitly do not specify the time argument of $\sigma_{IJ}$ in Eq~\eqref{eq:npi_tdc_final} because the \acrshort{npi} \acrshort{tdc} is not evaluated at a specific time but rather averaged over the entire time step. The time integral in Eq~\eqref{eq:npi_tdc_final} can be evaluated analytically, yielding a closed-form expression for the \acrshort{npi} \acrshort{tdc} and can be seen in Supporting Information of Ref.~\cite{10.1021/jz5009449}.

We can now start thinking about generalizing the \acrshort{npi} scheme to more than two states. First, let's differentiate overlap matrix $U_{IJ}$ in Eq~\eqref{eq:npi_unitary_matrix} with respect to $\Delta t$ to obtain
\begin{equation}\label{eq:npi_overlap_derivative}
\frac{\symrm{d}U_{IJ}}{\symrm{d}\Delta t}=\Braket{\phi_I(t)\vert\frac{\symrm{d}\phi_J(t+\Delta t)}{\symrm{d}\Delta t}}=\sum_K\Braket{\phi_I(t)\vert\phi_K(t+\Delta t)}\Braket{\phi_K(t+\Delta t)\vert\frac{\symrm{d}\phi_J(t+\Delta t)}{\symrm{d}\Delta t}},
\end{equation}
where on the right-hand side we have inserted a complete set of states at time $t+\Delta t$. We can now identify the last term in Eq~\eqref{eq:npi_overlap_derivative} as the \acrshort{tdc} at time $t+\Delta t$ and rearrange to obtain a matrix equation for the \acrshort{tdc} in terms of the overlap matrix and its derivative
\begin{equation}\label{eq:npi_matrix_equation}
\dot U=U\sigma\implies U(t+\Delta t)=e^{\int_t^{t+\Delta t}\sigma\symrm{d}t'}.
\end{equation}
Rearranging Eq~\eqref{eq:npi_matrix_equation} gives the formal solution for the multistate \acrshort{npi} \acrshort{tdc}
\begin{equation}\label{eq:npi_multistate_tdc}
\sigma=\frac{1}{\Delta t}\ln\left(U(t+\Delta t)\right).
\end{equation}
This expression reduces to the Eq~\eqref{eq:npi_tdc_final} in the two-state case. The matrix logarithm of a unitary matrix can be complex, so when implementing Eq~\eqref{eq:npi_multistate_tdc}, one must ensure that the diagonal elements of $U$ are positive to obtain a real-valued \acrshort{tdc}.
\subsection{The Baeck--An Scheme}\label{sec:baeck_an}
The approximations for the \acrshort{tdc} according to Baeck and An use, contrary to the previous methods, physical approximation. The original idea is to approximate the \acrshort{nacv} between two states around the conical intersection by a simple model. Baeck and An came up with the expression\supercite{10.1063/1.4975323}
\begin{equation}\label{eq:baeck_an_nacv}
\symbf{d}_{IJ}^{\symrm{BA}}=\frac{1}{2}\sqrt{\frac{\symrm{d}^2Z_{IJ}}{\symrm{d}R^2}\frac{1}{Z_{IJ}}},
\end{equation}
where $Z_{IJ}=\lvert E_J-E_I\rvert$ is the energy gap between adiabatic states $I$ and $J$. From this expression for the \acrshort{nacv}, we can directly obtain the \acrshort{tdc} using chain rule as
\begin{equation}\label{eq:lambda_tdc}
\sigma_{IJ}^{\symrm{\lambda}}=\frac{1}{2}\sqrt{\left(\ddot Z_{IJ}-\frac{\ddot R}{\dot R}\dot Z_{IJ}\right)\frac{1}{Z_{IJ}}},
\end{equation}
which is not a widely used expression for the \acrshort{tdc} so we have taken the liberty to denote it by $\lambda$ superscript, calling it the $\lambda$\acrshort{tdc}.\supercite{10.1021/acs.jctc.5c01176} The more popular expression is obtained by neglecting the second term in the parenthesis of Eq~\eqref{eq:lambda_tdc}, yielding
\begin{equation}\label{eq:kappa_tdc}
\sigma_{IJ}^{\kappa}=\frac{1}{2}\sqrt{\ddot Z_{IJ}\frac{1}{Z_{IJ}}},
\end{equation}
which is called the $\kappa$\acrshort{tdc}.\supercite{10.1021/acs.jctc.1c01080,10.1021/acs.jctc.3c00813,10.12688/openreseurope.13624.2} Both expressions in Eq~\eqref{eq:lambda_tdc} and Eq~\eqref{eq:kappa_tdc} have been shown to perform similarly in practice,\supercite{10.1021/acs.jctc.5c01176} with the $\kappa$\acrshort{tdc} being slightly more popular due to its simplicity. Both expressions diverge as $Z_{IJ}\to 0$, which is physically correct behavior for the \acrshort{tdc} near conical intersections. However, care must be taken when implementing these expressions to ensure numerical stability, especially the $\lambda$\acrshort{tdc} in Eq~\eqref{eq:lambda_tdc}, which involves division by the nuclear velocity $\dot R$.

These approximations have gained popularity in recent years due to their simplicity, but it has been shown that they can lead to significant errors in certain scenarios, especially when the assumptions underlying their derivation are not met. Therefore, while they can be useful tools in the computational chemist's toolkit, they should be applied with caution and validated against more accurate methods whenever possible.\supercite{10.1021/acs.jctc.5c01176}
\section{Ehrenfest Dynamics}\label{sec:ehrenfest_dynamics}
\section{Fewest Switches Surface Hopping}\label{sec:fewest_switches}
The \acrfull{fssh} method, developed by Tully in 1990,\supercite{10.1063/1.459170} is one of the most widely used \acrshort{tsh} algorithms. In \acrshort{fssh}, an ensemble of classical trajectories is propagated on adiabatic \acrshort{pes} according to Eq~\eqref{eq:tsh_eom}, with stochastic hops between electronic states determined by the instantaneous electronic populations. The key idea is to minimize the number of hops while ensuring that the ensemble-averaged populations match the quantum mechanical populations.

Befor we can obtain the hopping probabilities, it is convenient to express the Eq~\eqref{eq:tsh_eom} in terms of the density matrix elements $\rho_{IJ}=c_I c_J^\ast$ as
\begin{equation}\label{eq:fssh_density_matrix_eom}
i\dot\rho_{IJ}=\sum_K\left[\left(H_{IK}^{\symrm{el}}-i\sigma_{IK}\right)\rho_{KJ}-\left(H_{KJ}^{\symrm{el}}-i\sigma_{KJ}\right)\rho_{IK}\right],
\end{equation}
where the diagonal elements $\rho_{II}=\lvert c_I\rvert^2$ represent the populations of the adiabatic states, and the off-diagonal elements $\rho_{IJ}$ ($I\neq J$) represent the coherences between states. It is also useful to note that th electronic Hamiltonian (potential) is Hermitian and the \acrshort{tdc} is anti-Hermitian. From Eq~\eqref{eq:fssh_density_matrix_eom}, we can derive the time evolution of the populations $\rho_{II}$ using the properties of complex numbers as
\begin{equation}\label{eq:fssh_population_eom}
\dot\rho_{II}=\sum_K b_{KI},
\end{equation}
where
\begin{equation}\label{eq:fssh_b_matrix}
b_{KI}=2\Im\left(H_{IK}^{\symrm{el}}\rho_{IK}^\ast\right)+2\Re\left(\sigma_{IK}\rho_{IK}^\ast\right).
\end{equation}
If we consider an adiabatic basis where the electronic Hamiltonian is diagonal, the first term in Eq~\eqref{eq:fssh_b_matrix} vanishes, and the population transfer between states is mediated entirely by the \acrshort{tdc} term. We should also note that $b_{II}=0$ due to the properties of \acrshort{tdc} being anti-Hermitian and the density matrix being Hermitian.

Our goal is to derive hopping algorithm that minimizes the number of switches between states while ensuring that the ensemble-averaged populations match the quantum mechanical populations. Let's consider a two-state system for simplicity. In this case, the total number of trajectories at a time $t$ on state 1 be expressed as $N_1(t)=N\rho_{11}(t)$, where $N$ is the total number of trajectories in the ensemble. Without the loss of generality, we can assume that the population of state 1 is increasing, i.e., $\rho_{11}(t+\Delta t)>\rho_{11}(t)$. This means that the amount of trajectories that need to hop from state 2 to state 1 during the time interval $\Delta t$ must be larger than the number of trajectories hopping from state 1 to state 2. Since we want to minimize the number of hops, we can assume that no trajectories hop from state 1 to state 2 during this time interval. Therefore, the number of trajectories hopping from state 2 to state 1 must be equal to $N\left(\rho_{11}(t+\Delta t)-\rho_{11}(t)\right)$. The probability of a trajectory on state 2 hopping to state 1 during the time interval $\Delta t$ can then be expressed as the ratio of the number of trajectories hopping to the total number of trajectories on state 2 as
\begin{equation}\label{eq:fssh_hopping_probability_two_state}
    P_{2\to 1}=\frac{N\left(\rho_{11}(t+\Delta t)-\rho_{11}(t)\right)}{N\rho_{22}}=\frac{\dot\rho_{11}\Delta t}{\rho_{22}}=\frac{b_{12}\Delta t}{\rho_{22}},
\end{equation}
where we used $\dot\rho_{11}=b_{12}$ from Eq~\eqref{eq:fssh_population_eom}. The ho probability from state 1 to state 2 can be derived similarly assuming that the population of state 1 is decreasing during the time interval $\Delta t$.

Should we want to generalize the hopping probabilities to more than two states, we can follow a similar reasoning. Let's assume that the population of state $I$ is increasing during the time interval $\Delta t$. The number of trajectories hopping from all other states to state $I$ must then be equal to $N\left(\rho_{II}(t+\Delta t)-\rho_{II}(t)\right)$. To minimize the number of hops, we can assume that no trajectories hop from state $I$ to any other state during this time interval. Therefore, the probability of a trajectory on state $J$ hopping to state $I$ during the time interval $\Delta t$ can be expressed as
\begin{equation}\label{eq:fssh_hopping_probability_multistate}
P_{J\to I}=\frac{b_{IJ}\Delta t}{\rho_{JJ}},
\end{equation}
where in the nominator we have the population flux only from state $J$ to state $I$, ignoring contributions from other states.
\section{Landau--Zener Surface Hopping}\label{sec:landau_zener}
\section{Mapping Approach to Surface Hopping}\label{sec:mapping_approach}
\section{Beyond Surface Hopping}\label{sec:beyond_surface_hopping}
